\section{Excluding one complete affine extension level}\label{sec:affine}

The product construction below belongs to the extension-Nullstellensatz
approach developed by Buss, Impagliazzo, Kraj\'i\v{c}ek, Pudl\'ak,
Razborov, and Sgall \cite{BIKPRS}. We use a complete one-level family
and give the needed algebraic arguments explicitly.

An accuracy-\(h\) affine block has a finite tuple of old affine polynomials
\(g_1,\ldots,g_s\), fresh coefficient variables \(r_{ui}\), and product
\begin{equation}\label{eq:block}
 P=\prod_{u=1}^h\left(1-\sum_{i=1}^s r_{ui}g_i\right).
\end{equation}
Its \emph{complete} axiom set contains every \(g_iP\) and every
\(r_{ui}^2-r_{ui}\). Products \(P=0\) are not added as axioms.
Different blocks have disjoint coefficient variables, and every input
is affine in the old bits alone.

\subsection{Cleanup and ordinary affine restriction}
A block is \emph{proper} if its affine input span is nonzero and does
not contain the constant one. Its rank is the dimension of that span.

\begin{lemma}[Scalar cleanup]\label{lem:cleanup}
Zero-span and unit-span blocks can be removed by substituting constants
for their coefficients, without increasing PC degree or changing the
other blocks.
After zero inputs are discarded, every companion in a proper block
has ordinary degree \(2h+1\).
\end{lemma}
\begin{proof}
For zero span all inputs are zero; set its coefficients to zero.
For unit span choose constants \(c_i\) with \(\sum_i c_i g_i=1\).
Set the first coefficient row to \(c_i\) and all other rows to zero,
making \(P=0\).
Every axiom of an eliminated block maps to zero, and coefficient domains
map to valid constants. Other blocks use only old bits and are unchanged.
A constant substitution preserves a PC derivation and its degree.

In a proper span each nonzero input has nonzero linear part; otherwise
the span contains one. Each factor of \eqref{eq:block} then has nonzero
degree-two part, because its own coefficient variables are distinct.
Their product has degree \(2h\) in the ordinary polynomial ring.
Multiplying by a nonzero affine input gives degree \(2h+1\).
Coefficients for an identically zero input may independently be set to
zero, so those inputs cause no exceptional axiom degree.
\end{proof}

For a proper block, an affine basis \(F_1,\ldots,F_r\) of its input span
has independent linear parts. Indeed, a nontrivial relation among those
parts would give either a dependence among the \(F_j\) or a nonzero
constant in their span. Thus its zero set is a nonempty affine flat of
codimension \(r\).

\begin{lemma}[A common ordinary restriction kernel]\label{lem:kernel}
Let \(h=3\ell\), \(m=n+1\), \(v=m\ell\), and let \(M\ge1\) bound the
number of proper blocks. Put
\[
 k=\left\lceil\sqrt{m\ln(4M)}\right\rceil
\]
and suppose \(k\le m\).
For all blocks of rank greater than \(h(k+1)\), there is a common nonzero
\(f\in\mathcal L_k\) whose restriction to each zero flat is the zero
\emph{ordinary polynomial} in affine free coordinates.
For each such block it has a literal representation
\begin{equation}\label{eq:literal-ideal}
 f=\sum_i a_i g_i,\qquad \deg a_i\le k-1.
\end{equation}
\end{lemma}
\begin{proof}
Restriction to a codimension-\(r\) affine flat maps \(\mathcal L_k\)
into ordinary polynomials of degree at most \(k\) on \(v-r\) free
coordinates. Its image dimension is at most
\(\binom{v-r+k}{k}\).
Every high-rank block has \(r\ge r_*=3\ell(k+1)+1\).
If \(r_*>v\), there are no such blocks and any nonzero \(f\) suffices.
Otherwise
\begin{align}
 \frac{\binom{v-r_*+k}{k}}{\binom mk\ell^k}
  &=\prod_{j=0}^{k-1}
      \left(1-\frac{r_*-1-(\ell+1)j}{v-\ell j}\right)
       \label{eq:kernel-ratio}\\
  &\le \left(1-\frac{k}{m}\right)^k
   \le e^{-k^2/m}. \nonumber
\end{align}
For the equality, order the numerator factors increasingly, as
\(v-r_*+1+j\), and the denominator factors as \(v-\ell j\).
The numerator of each subtracted fraction is at least \(\ell k\);
its denominator is at most \(m\ell\).
The presence of a high block implies \(k<m\), so all factors used here
are nonnegative.

The sum of all high restriction-image dimensions is consequently at
most
\[
 M e^{-k^2/m}\binom mk\ell^k
   \le\frac14\binom mk\ell^k
   <\dim\mathcal L_k.
\]
The joint restriction map has a nonzero kernel.

For a fixed high block, complete \(F_1,\ldots,F_r\) to affine coordinates.
The ordinary polynomial for \(f\) has zero specialization when the first
\(r\) coordinates are zero. Every monomial therefore has one of these
coordinates as a factor. Assign each such monomial to one of its factors,
obtaining \(f=\sum_j F_j A_j\) with \(\deg A_j\le k-1\).
Pull back the affine coordinate change and express each \(F_j\) as a
constant linear combination of the original \(g_i\).
This gives \eqref{eq:literal-ideal} at the same degree.
\end{proof}

\subsection{One simultaneous weighted substitution}
\begin{lemma}[Weighted removal of the affine family]\label{lem:hybrid}
Let an old system \(\mathcal F\) contain all old Boolean equations,
and adjoin a complete proper one-level affine family of accuracy \(h\).
Suppose it has an ordinary-PC refutation of degree \(D\ge2h+1\).
Let \(k\ge1\), and let \(f\) be a multilinear polynomial of degree at most
\(k\) having representations \eqref{eq:literal-ideal} for every block
of rank greater than \(h(k+1)\).
Then
\[
 f\in\Cons_{k(D+1)}(\mathcal F).
\]
\end{lemma}
\begin{proof}
We give one substitution \(\sigma\) fixing the old variables and sending
every coefficient variable to an old polynomial of degree at most \(k\).

\paragraph{Low rank.}
For a block of rank \(r\le h(k+1)\), choose an affine basis
\(F_1,\ldots,F_r\) and partition it into \(h\) ordered groups of size at
most \(k+1\), allowing empty groups. In a group \(J\), telescoping gives
\[
 1-\prod_{j\in J}(1-F_j)
   =\sum_{j\in J}F_j\prod_{\substack{t\in J\\t<j}}(1-F_t).
\]
Express the \(F_j\) in the original \(g_i\). The coefficient of each
\(g_i\) on the right has degree at most \(k\); assign it to the
corresponding row of \eqref{eq:block}.
The block product becomes
\[
 Z=\prod_{j=1}^r(1-F_j).
\]
At every Boolean point, \(Z\) is the indicator that all inputs vanish.
Thus \(g_iZ\) is a zero Boolean function of ordinary degree at most
\(r+1\). Lemma~\ref{lem:boolean} gives its certificate through \(r+1\).
After multiplication by \(f\), each companion image is certified through
\begin{equation}\label{eq:low-cost}
 r+1+k\le h(k+1)+1+k\le k(2h+2).
\end{equation}
The last inequality is equivalent to \((k-1)(h+1)\ge0\).

\paragraph{High rank.}
Use \eqref{eq:literal-ideal} as the first coefficient row and set the
other rows to zero. The product becomes \(1-f\). The weighted companion
image is
\[
 f g_i(1-f)=-g_i(f^2-f).
\]
By Lemma~\ref{lem:boolean}, it has an old Boolean certificate through
\(2k+1\le k(2h+2)\).

\paragraph{All remaining axioms.}
If a coefficient variable maps to \(\beta\), its weighted field equation
is \(f(\beta^2-\beta)\), certified through \(3k\).
An old degree-\(e\) axiom \(F\) has weighted image \(fF\), with certificate
through \(e+k\le k(e+1)\).
Together with \eqref{eq:low-cost}, this yields the complete ledger
\begin{center}
\begin{tabular}{lll}
\toprule
Original axiom & Original degree \(e\) & Weighted image ceiling\\
\midrule
High companion & \(2h+1\) & \(2k+1\)\\
Low companion & \(2h+1\) & \(r+1+k\)\\
Coefficient Booleanity & \(2\) & \(3k\)\\
Old axiom & \(e\) & \(e+k\)\\
\bottomrule
\end{tabular}
\end{center}
Every entry fits \(k(e+1)\). Proper-block companion degrees are the
ordinary degrees justified in Lemma~\ref{lem:cleanup}.
All assignments use only old variables, so they define a single
substitution even when input tuples overlap.

\paragraph{Primitive PC replay.}
Let \(B=k(D+1)\). Replace each source line \(q\) by \(f\sigma(q)\).
The preceding certificates initialize every source axiom used by the
degree-\(D\) refutation through at most \(B\).
Linear combinations commute with the weighted map.
For a nonzero source multiplication \(q\mapsto yq\), write \(d=\deg q\).
Then \(d+1\le D\), while
\[
 \deg(f\sigma(q))\le k(d+1),\qquad \deg\sigma(y)\le k.
\]
Multiplying the completed weighted line by \(\sigma(y)\), using
Lemma~\ref{lem:pc-product}, therefore stays within
\(k(d+2)\le B\). Zero source lines cause no difficulty.
The final weighted image of one is \(f\), proving the assertion.
\end{proof}

\begin{theorem}[Finite-parameter affine exclusion]\label{thm:affine}
Let \(n=2^\ell\), \(\ell\ge2\), \(m=n+1\), and \(h=3\ell\).
Adjoin any finite complete one-level old-affine family to \(\mathcal Q_n\).
Let \(M\ge1\) bound its number of proper blocks after scalar cleanup,
and put
\[
 k=\left\lceil\sqrt{m\ln(4M)}\right\rceil,\qquad B=k(D+1).
\]
If
\begin{equation}\label{eq:affine-room}
 D\ge2h+1,\qquad k\le m,\qquad n\ge2B-1,\qquad4(k-1)<n,
\end{equation}
then the augmented system has no ordinary-PC refutation through degree \(D\).
In particular this excludes polynomially many blocks at every
polylogarithmic degree for sufficiently large \(n\).
\end{theorem}
\begin{proof}
Apply scalar cleanup, which preserves a hypothetical refutation and
its degree. Lemma~\ref{lem:kernel} supplies a nonzero
\(f\in\mathcal L_k\) with the representations required by
Lemma~\ref{lem:hybrid}. The latter gives
\(f\in\Cons_B(\mathcal Q_n)\).
Since \(k\le B\), equation~\eqref{eq:affine-room} contradicts
Theorem~\ref{thm:row-separation}. If no proper blocks remain, the same
argument may use \(f=1\), or directly use the normalized old functional.

For polynomial \(M\), \(k=O(\sqrt{n\log n})\).
For any fixed polylogarithmic bound on \(D\),
\(B=\sqrt n\,\operatorname{polylog}n=o(n)\), so the remaining conditions
hold eventually. If the proposed degree bound is smaller than \(2h+1\),
enlarge it to \(2h+1\); this remains polylogarithmic and still excludes
the original refutation.
\end{proof}
